% study-pack skeleton: FORMULA SHEET. Copy to src/, rename, replace every <...>. Keep the structure.
\documentclass[11pt,a4paper]{article}
\usepackage[margin=1.8cm,headheight=14pt]{geometry}
\def\studylang{english}
\def\sourcelang{english}
\input{preamble}
\studyheader{<Course · Chapter>}{<Formula sheet>}
\begin{document}
{\LARGE\bfseries\sffamily\color{accent} <Title>}\par\smallskip
\begin{center}
\begin{tikzpicture}[node distance=8mm and 14mm,every node/.style={draw=accent,rounded corners,font=\footnotesize,inner sep=3pt}]
\node (a) {(1) <f>}; \node (b) [right=of a] {(2) <f'>}; \node (c) [right=of b] {(3) <g>};
\draw[-Stealth] (a) -- node[draw=none,above,font=\scriptsize]{<derivative>} (b);
\draw[-Stealth] (b) -- node[draw=none,above,font=\scriptsize]{<substitution>} (c);
\end{tikzpicture}
\end{center}
\renewcommand{\arraystretch}{1.25}
\begin{tabularx}{\linewidth}{r>{\raggedright\arraybackslash}X>{\raggedright\arraybackslash}p{5.2cm}}\toprule
\# & Formula & From what, how\\\midrule
1 & $<formula>$ & definition\\
2 & $<formula>$ & from (1): <derivative>\\
3 & $<formula>$ \hfill\textit{bonus} & from (2): <substitution>\\\bottomrule
\end{tabularx}
\begin{method}[: <typical problem>]
\begin{enumerate}\item <step> \item <step>\end{enumerate}
\end{method}
\end{document}
